\documentclass[hidelinks, 11 pt]{exam}

\usepackage[a4paper, top=25.4mm, bottom=25.4mm, left=19.1mm, right=19.1mm, footskip=8mm]{geometry}
\usepackage{parskip}
\usepackage[list=true, listformat=simple]{subcaption}
\usepackage{graphicx}
\usepackage{float}
\usepackage{amsthm}
\usepackage{mathtools}
\usepackage{amsfonts}
\usepackage[dvipsnames]{xcolor}
\usepackage{tikz}
\usepackage{pgfplots}
\usepackage{tkz-euclide}
\usepackage{multirow}
\usepackage{hyperref}
\usepackage{titlesec}
\usepackage{array}
\usepackage{booktabs}
\usepackage{enumitem}
\usepackage{environ}
\usepackage{etoolbox}
\usepackage{needspace}
\usepackage{footnote}

\usetikzlibrary{
	quotes,
	angles,
	positioning,
	math,
	calc,
	decorations.markings,
	arrows.meta}

\widowpenalty=10000
\clubpenalty=10000

\newcommand{\tocsection}[1]{%
	\refstepcounter{section}%
	\addcontentsline{toc}{section}{\protect\numberline{\thesection}#1}%
}

% Redefining space around figures and tables
\AtBeginEnvironment{figure}{\vspace{-1mm}}
\AtEndEnvironment{table}{\vspace{-4mm}}

% Redefining space around figure captions
\captionsetup[figure]{aboveskip=2mm, belowskip=-2mm}
\captionsetup[subfigure]{aboveskip=2mm, belowskip=1mm}

% Spacing around headings
\titlespacing*{\section}{0pt}{10pt}{5pt}
\titlespacing*{\subsection}{0pt}{5pt}{5pt}

% Other spacing
\setlength{\parskip}{7pt plus 4pt}
\setlength{\parindent}{0mm} % Start of paragraph line indent
\setlength{\skip\footins}{5mm} % Footnotes

\makeatletter
\renewcommand\paragraph{\@startsection{paragraph}{4}{\z@}%
	{-2.5ex\@plus -1ex \@minus -0.25ex}%
	{1.25ex \@plus 0.25ex}%
	{\normalfont\normalsize\bfseries}}
\makeatother
\let\oldthesubsection\thesubsection

\hypersetup{
	colorlinks = true, % Colours links instead of ugly boxes
	urlcolor = blue,
	linkcolor = blue,
	citecolor = blue,
}

% Dark mode

\definecolor{dark}{HTML}{363636}
\definecolor{light}{HTML}{f5f5f5}

\colorlet{myblack}{dark}
\colorlet{mywhite}{light}

\newcommand{\StyleDark}{
	\colorlet{myblack}{light}
	\colorlet{mywhite}{dark}
	\pagecolor{mywhite}
	\color{myblack}
	\colorlet{mylinkblue}{blue!10}
	
	\hypersetup{
		colorlinks=true,
		urlcolor=mylinkblue,
		linkcolor=mylinkblue,
		citecolor=mylinkblue,
		filecolor=mylinkblue,
	}
}

% Make hints environment.
\makeatletter

\newif\ifhashints
\hashintsfalse

\newcommand{\storedhints}{}

\NewEnviron{hints}{%
	\global\hashintstrue
	\xdef\storedhints{\unexpanded\expandafter{\BODY}}%
}

\let\origsolutiontitle\solutiontitle

\renewcommand{\solutiontitle}{%
	\ifhashints
	\noindent\textbf{Hints:}\par
	\storedhints
	\medskip\par
	\fi
	\origsolutiontitle
}

\AfterEndEnvironment{solution}{%
	\global\hashintsfalse
	\global\let\storedhints\empty
}

\makeatother

% Make footnotes work in dynamically shown environements.
\makesavenoteenv{hints}
\makesavenoteenv{solution}

% Define arrow on part of the line
% For example \path[->-=0.5] (A) -- (B)
\tikzset{->-/.style={decoration={
			markings,
			mark=at position #1 with {\arrow[xshift=3]{Straight Barb [length=7, width=7]}}}, postaction={decorate}}}

\tikzset{
	>= {Stealth[scale=1.5]}
}

% Label the length of a line
% \for example the below draws a distance label of 10 that is 0.2 above the line AB.
% \pic {dimension={$(A)$}{$(B)$}{0.2}{$10$}};
\tikzset{
	pics/.cd,
	dimension/.style args={#1#2#3#4}{
		code={
			\draw[|<->|] ($(#1)!#3!-90:(#2)$) -- node[midway, fill=mywhite, blend mode=normal] {#4} ($(#2)!#3!90:(#1)$);
		}
	}
}